% 第 15 题：山体截面与机器人表演装置（按原图矢量重绘）
\begin{tikzpicture}[font=\small, >=Stealth, scale=0.95]
  % 山体轮廓
  \fill[gray!12] (8.2,7.5) -- (10.0,7.5) -- (10.0,0.2) -- (7.4,0.2)
      .. controls (7.0,1.2) and (7.2,2.2) .. (6.4,3.0)
      .. controls (5.9,3.6) and (6.0,4.3) .. (5.6,5.1)
      .. controls (5.2,5.9) and (5.6,6.5) .. (6.2,7.0)
      .. controls (7.0,7.4) and (7.6,7.5) .. (8.2,7.5) -- cycle;
  \draw[thick] (8.2,7.5) -- (10.0,7.5) -- (10.0,0.2) -- (7.4,0.2)
      .. controls (7.0,1.2) and (7.2,2.2) .. (6.4,3.0)
      .. controls (5.9,3.6) and (6.0,4.3) .. (5.6,5.1)
      .. controls (5.2,5.9) and (5.6,6.5) .. (6.2,7.0)
      .. controls (7.0,7.4) and (7.6,7.5) .. (8.2,7.5);
  % A 点（崖顶）
  \fill (6.35,6.95) circle (1.8pt) node[above left] {$A$};
  % 两条平行轨道 a、b
  \draw[thick] (0,4.55) -- (6.55,4.55);
  \draw[thick] (0,3.75) -- (6.55,3.75);
  \node[left] at (-0.1,4.55) {$a$};
  \node[left] at (-0.1,3.75) {$b$};
  % 滑杆 M（与轨道垂直的短斜线示意）
  \draw[line width=2.2pt] (2.35,3.55) -- (2.75,4.75);
  \node[above left] at (2.4,4.7) {$M$};
  % 机器人（B 点处，下方小圆）
  \draw (3.95,3.35) circle (0.11);
  \node[right] at (4.1,3.3) {$m$};
  % 轻绳 M--机器人
  \draw[dashed] (2.65,4.55) -- (3.95,3.35);
  % B 点与向下初速度
  \fill (3.95,3.75) circle (1.5pt);
  \node[above left] at (3.9,3.8) {$B$};
  \draw[->,thick] (3.95,3.2) -- (3.95,2.4) node[below] {$v$};
  % 高度标注
  \draw[<->] (3.35,3.75) -- (3.35,6.95);
  \node[left] at (3.35,5.35) {$1.2L$};
  \draw[<->] (0.55,3.75) -- (0.55,4.55);
  \draw[<->] (2.05,3.75) -- (3.95,3.75);
  \node[below] at (3.0,3.75) {$L$};
  % 轨道延长虚线（滑杆位置到左端）
  \draw[dashed] (0.55,3.1) -- (0.55,3.75);
  \draw (0.55,3.0) circle (0.07);
\end{tikzpicture}
